
\documentclass[12pt]{article}

\usepackage[margin=1in]{geometry}
\usepackage{setspace}
\usepackage{times}
\usepackage{indentfirst}
\usepackage{hyperref}

\onehalfspacing

\begin{document}

\begin{center}
\textbf{Radar Operations on the Day of the Pearl Harbor Attack} \\[0.5em]
April 27, 2025 \\[0.5em]
1480 words
\end{center}

\vspace{1em}

The salty ocean breeze blew through the windows of the operating van and onto Private George Elliott’s cheeks. Looking out from Opana Point, a location 532 feet above sea level, there was an unobstructed view of the Pacific waters. Even in December, it was another welcoming, cloudless day at the tropical home of the U.S. Pacific Fleet.

But right now, it was not the time for Elliott to get distracted. Joseph Lockard, the more experienced radar operator, was teaching Elliott how to use the radar oscilloscope. Two years ago, six mobile sets of this long-range aircraft search apparatus called the SCR-270 radar had been installed across Hawaii. Each set consisted of four trucks that carried a transmitter, modulator, water cooler, receiver, oscilloscope, operator, and generator. Along with that came a 45-foot-tall rotating metal antenna that emits and receives radio waves. This device could detect planes and ships before the naked eye could see them. On the morning of December 7, 1941, the Opana Radar Station was the only radar station in operation. Higher-ranking officers had limited the use of the radar sets, fearing that the equipment would wear out.

The attack had begun.

\begin{thebibliography}{9}

\bibitem{ref1}
Commander, Navy Region Hawaii,
\textit{History},
28 February 2025.
Available: \url{https://cnrh.cnic.navy.mil/About/History/}

\end{thebibliography}

\end{document}
